\documentclass[12pt, a4paper]{article}
\usepackage[T1]{fontenc}
\usepackage[utf8]{inputenc}
\usepackage[spanish]{babel}
\usepackage{geometry}
\usepackage{graphicx}
\usepackage{xcolor}
\usepackage{booktabs}
\usepackage{hyperref}
\usepackage{palatino} % Fuente más amigable para lectura
\usepackage{amsmath} % Para fórmulas matemáticas

\geometry{top=2.5cm, bottom=2.5cm, left=2.5cm, right=2.5cm}

\hypersetup{
    colorlinks=true,
    linkcolor=blue,
    filecolor=magenta,      
    urlcolor=cyan,
}

\title{\textbf{Análisis de Sistemas en Frecuencia}\\ \large Entendiendo los Diagramas de Bode}
\author{Generado por Gemini}
\date{\today}

\begin{document}

\maketitle
\tableofcontents
\newpage

\section{Introducción: ¿Qué es un Diagrama de Bode?}

Un \textbf{Diagrama de Bode} (pronunciado 'Bodi') es una herramienta gráfica fundamental en ingeniería de control, electrónica y procesamiento de señales. Su propósito es mostrar cómo responde un sistema lineal e invariante en el tiempo (LTI) a diferentes frecuencias.

En lugar de ser un solo gráfico, un diagrama de Bode consta de \textbf{dos gráficos} que se analizan en conjunto:
\begin{enumerate}
    \item \textbf{El Diagrama de Magnitud:} Muestra la ganancia (cuánto se amplifica o atenúa la señal) en función de la frecuencia.
    \item \textbf{El Diagrama de Fase:} Muestra el desfase (cuánto se retrasa o adelanta la señal) en función de la frecuencia.
\end{enumerate}

La clave de su utilidad es que ambos gráficos utilizan una \textbf{escala logarítmica} para el eje de la frecuencia, lo que permite analizar un rango muy amplio de comportamiento, desde frecuencias muy bajas hasta muy altas, en un solo vistazo.

\section{Los Dos Componentes del Diagrama}

\subsection{El Diagrama de Magnitud (Ganancia)}
Este gráfico nos dice cómo cambia la "fuerza" de la señal al pasar por el sistema.
\begin{itemize}
    \item \textbf{Eje Vertical (Y):} La magnitud se mide en \textbf{decibelios (dB)}. La fórmula para convertir la ganancia $G$ a decibelios es:
    $$ G_{dB} = 20 \log_{10}(|H(j\omega)|) $$
    Donde $|H(j\omega)|$ es la magnitud de la función de transferencia del sistema.
    \begin{itemize}
        \item Una ganancia de \textbf{0 dB} significa que la señal no se amplifica ni se atenúa.
        \item Una ganancia \textbf{positiva} (ej. +6 dB) significa amplificación.
        \item Una ganancia \textbf{negativa} (ej. -3 dB) significa atenuación.
    \end{itemize}
    \item \textbf{Eje Horizontal (X):} La frecuencia ($\omega$ en rad/s o $f$ en Hz) se representa en escala logarítmica.
\end{itemize}

\subsection{El Diagrama de Fase}
Este gráfico nos dice cómo se desplaza en el tiempo la señal de salida con respecto a la de entrada.
\begin{itemize}
    \item \textbf{Eje Vertical (Y):} El desfase se mide en \textbf{grados} o radianes.
    \begin{itemize}
        \item Una fase de \textbf{0°} significa que la salida está perfectamente sincronizada con la entrada.
        \item Una fase \textbf{negativa} (ej. -90°) indica que la salida está retrasada con respecto a la entrada.
        \item Una fase \textbf{positiva} indica un adelanto.
    \end{itemize}
    \item \textbf{Eje Horizontal (X):} La frecuencia, también en escala logarítmica.
\end{itemize}

\section{¿Para Qué Sirven los Diagramas de Bode?}
Su popularidad se debe a que simplifican enormemente el análisis de sistemas complejos. Sus usos principales son:

\begin{itemize}
    \item \textbf{Análisis de Estabilidad:} En sistemas realimentados (como un amplificador con control de volumen o un piloto automático), los diagramas de Bode permiten calcular el \textbf{margen de ganancia} y el \textbf{margen de fase}. Estos dos valores nos dicen qué tan lejos está el sistema de volverse inestable y empezar a oscilar sin control.
    
    \item \textbf{Diseño de Filtros:} Son la herramienta por excelencia para visualizar el comportamiento de filtros electrónicos. Es muy fácil ver si un circuito se comporta como un filtro \textbf{pasa-bajo} (deja pasar solo frecuencias bajas), \textbf{pasa-alto} (deja pasar altas) o \textbf{pasa-banda} (deja pasar un rango específico).
    
    \item \textbf{Caracterización de Sistemas:} Permiten identificar la \textbf{frecuencia de corte} (donde el sistema empieza a atenuar), el \textbf{ancho de banda} y los picos de \textbf{resonancia} de cualquier sistema, desde un altavoz hasta la suspensión de un coche.
\end{itemize}

\section{Ejemplo Conceptual: Un Filtro Pasa-Bajo Simple}
Imaginemos un circuito simple que solo deja pasar los sonidos graves (bajas frecuencias) y bloquea los agudos (altas frecuencias). Su diagrama de Bode se vería así:

\begin{itemize}
    \item \textbf{Diagrama de Magnitud:}
    \begin{itemize}
        \item A frecuencias bajas, la línea sería horizontal y estaría en 0 dB. Esto significa que los sonidos graves pasan sin problemas.
        \item A partir de una "frecuencia de corte", la línea empezaría a caer con una pendiente constante (típicamente -20 dB por década). Esto muestra cómo los sonidos agudos son cada vez más atenuados.
    \end{itemize}
    \item \textbf{Diagrama de Fase:}
    \begin{itemize}
        \item A frecuencias muy bajas, la fase sería 0°.
        \item A medida que la frecuencia aumenta, la fase iría cayendo hasta estabilizarse en -90°, mostrando que las señales más altas sufren un mayor retardo.
    \end{itemize}
\end{itemize}

\section{Conclusión}
En resumen, los diagramas de Bode son el "electrocardiograma" de un sistema electrónico o de control. Proporcionan una visión completa y rápida de cómo un sistema interactuará con las señales del mundo real, que están compuestas por una infinidad de frecuencias. Son una herramienta indispensable para cualquier ingeniero que trabaje con sistemas dinámicos.

\end{document}